\section{Whole actual normal-series currying and relative polynomial flatness}

\begin{theorem}[Whole formal normal coefficients give an actual algebra equivalence]
\label{hir:actual-whole-formal-normal-series-equivalence}
\uses{hir:actual-formal-centre-power-relative-coefficients}
\lean{Hironaka.HironakaActualFormalNormalSeries.formalNormalSeriesEquiv,Hironaka.HironakaActualFormalNormalSeries.formalNormalSeriesEquiv_coefficient}
\leanok
Literal entire normal coefficient extraction gives a genuine algebra
isomorphism
\[
 \mathbb C[[u_1,\ldots,u_d,z_1,\ldots,z_t]]
 \simeq_{\mathbb C}
 \mathbb C[[u_1,\ldots,u_d]][[z_1,\ldots,z_t]],
 \quad F\longmapsto(\gamma\mapsto F_\gamma).
\]
Every coefficient is the original coefficient at the joined exponent.
This is an entire series-ring equivalence, including empty coordinate
blocks, without a formal completion identification supplied as input.
\end{theorem}
\begin{proof}
Split each actual exponent into its literal first and second coordinate
blocks. Their joined exponent recovers it, and splitting a joined
exponent recovers both blocks. Consequently coefficient currying and
uncurrying are inverse functions on the entire series spaces.
Zero, addition and constants are retained coefficientwise. For
multiplication, a decomposition of a joined exponent into two
nonnegative exponents is equivalent to one decomposition of its
normal block and one of its tangent block. The splitting and joining
identities give an actual bijection between these finite antidiagonal
index types. Reindex the genuine convolution sum by this bijection;
the two finite convolutions are exactly the product in the iterated
power-series ring. This proves multiplicativity on whole series.
The same literal coefficient formula proves preservation of complex
constants, giving the stated algebra equivalence.
\end{proof}

\begin{theorem}[Actual relative polynomial evaluation has its literal whole Taylor scalar dictionary]
\label{hir:actual-relative-polynomial-whole-taylor-dictionary}
\uses{hir:actual-whole-formal-normal-series-equivalence,hir:actual-relative-polynomial-formal-taylor-coefficient}
\lean{Hironaka.HironakaActualFormalNormalSeries.formalNormalSeries_normalPolynomialEval}
\leanok
Put $B=\mathcal O_{\mathbb C^d,0}$ and
$\widehat B=\mathbb C[[u_1,\ldots,u_d]]$. For every literal
$P\in B[z_1,\ldots,z_t]$, curry the actual whole Taylor series of its
actual analytic normal-polynomial evaluation. The result is exactly
$P$ with each coefficient replaced by its actual whole Taylor series,
then included into $\widehat B[[z]]$. The formal first-block scalar
embedding is therefore the actual Taylor image of the genuine analytic
first-block pullback.
\end{theorem}
\begin{proof}
For a constant polynomial with actual analytic coefficient $b$, the
genuine whole Taylor naturality for the tangent inclusion identifies
its Taylor series with the first-block rename of the actual Taylor
series of $b$. Literal joined coefficient extraction of this rename
is zero at nonzero normal exponents and the same original series at
normal exponent zero. This proves the constant case. For each variable,
actual analytic normal evaluation is the genuine normal coordinate
germ, whose actual Taylor series is its literal formal variable.
Currying its coefficients gives the corresponding iterated formal
normal variable. Both sides are ring homomorphisms, so polynomial
induction by constants, addition and multiplication by a variable
proves the whole equality for every polynomial. The constant case
also proves the entire scalar compatibility.
\end{proof}

\begin{theorem}[Entire actual relative polynomial evaluation is flat]
\label{hir:actual-relative-polynomial-analytic-flatness}
\uses{hir:actual-relative-polynomial-whole-taylor-dictionary,hir:actual-original-polynomial-analytic-flatness,hir:actual-original-taylor-faithfully-flat}
\lean{Hironaka.HironakaActualFormalNormalSeries.normalPolynomialEval_flat}
\leanok
The genuine entire ring map
$B[z_1,\ldots,z_t]\to\mathcal O_{\mathbb C^{d+t},0}$
given by original tangent pullback on every analytic coefficient and
literal original normal coordinate evaluation is flat. No relative
flatness, desired completion equality or analytic relation-generation
certificate is supplied.
\end{theorem}
\begin{proof}
The actual Taylor map $B\to\widehat B$ is faithfully flat, in
particular flat. Its genuine polynomial base change is flat: the
actual tensor-product polynomial equivalence identifies
$\widehat B\otimes_B B[z]$ with $\widehat B[z]$ and retains each
literal polynomial coefficient map. The Noetherian polynomial ring
$\widehat B[z]$ has flat actual completion in the ideal of its normal
variables. The genuine polynomial variables-adic completion
isomorphism identifies this completion with $\widehat B[[z]]$ over
that whole polynomial ring. Thus the composite map from $B[z]$ to
$\widehat B[[z]]$ is flat. The first theorem transports it to the
whole formal ring in all actual coordinates. The second theorem
identifies the resulting ring map with actual relative polynomial
evaluation followed by the actual whole Taylor map of
$\mathcal O_{\mathbb C^{d+t},0}$. That latter Taylor map is
faithfully flat. Apply the genuine faithful-target flatness descent
proved earlier, with the exact scalar tower just established. The
result is flatness of the stated entire original analytic evaluation.
\end{proof}

\begin{corollary}[Every actual analytic relation on relative polynomials is generated by whole polynomial syzygies]
\label{hir:actual-relative-polynomial-analytic-relations}
\uses{hir:actual-relative-polynomial-analytic-flatness}
\lean{Hironaka.HironakaActualFormalNormalSeries.relative_analytic_relation_iff_polynomial_syzygies}
\leanok
For any finite family $p_i\in B[z]$ and actual analytic germs
$A_i\in\mathcal O_{\mathbb C^{d+t},0}$, the literal whole identity
$\sum_i\operatorname{ev}(p_i)A_i=0$ holds if and only if there are
finitely many whole polynomial syzygies $c_{il}\in B[z]$ and actual
analytic germs $b_l$ such that
$A_i=\sum_l\operatorname{ev}(c_{il})b_l$ and
$\sum_i p_ic_{il}=0$. The finite number, polynomial coefficients
and actual analytic multipliers are all produced internally.
\end{corollary}
\begin{proof}
Equip the actual target analytic ring with the scalar action of the
literal relative polynomial evaluation map. The proved whole flatness
and the genuine finite equational criterion for flat modules express
every finite relation as a finite sum of whole relations over that
polynomial ring. Under the actual scalar-action formula these identities
are precisely the displayed actual analytic multiplier identities.
Conversely substitute the finite multiplier expressions, distribute
multiplication and interchange the finite sums. Every polynomial column
has zero entire sum, and actual evaluation preserves that zero, so the
original analytic sum is zero. No output syzygy generation is assumed.
\end{proof}

No exceptional relation lifting or relative normalized saturation is
asserted here. These require their own genuine coefficient and syzygy
arguments after this entire relative flatness producer.
